\section{Discussion}
\subsection{Accuracy and the Variational Principle}
\label{sec:hwe_limitation}
The most significant finding within this stack is the relationship between the ansatz expressibility and physical 
correctness. The Hardware Efficient (HWE) Ansatz was a practical choice for NISQ hardware due to its shallow circuit depth 
allowing circuits to complete before qubits lose coherency~\parencite{kandala2017hardware}, but is limited as an ansatz in conserving particle number. 
During the optimization stages, the SPSA optimizer can drive variational parameters into regions of the Hilbert space
where the prepared state $|\psi(\theta)\rangle$ contains superpositions of different electron numbers, creating states that are
unphysical in molecular chemistry. This results in the energy of the VQE dropping below the established, exact FCI ground state energy, 
which violates the variational principle within the correct particle number sector. Consider the molecule $H_{2}$, where the  
optimizer crossed below FCI at approximately iteration 126. For the larger molecule $H_{2}O$, this occurred as early as iteration 37. The rate 
at which this violation occurs has been found to correlate directly with the ratio of qubits to variational parameters. The molecule 
$H_{2}O$ (112 parameters, 14 qubits) explores the parameter space more quickly than $H_{2}$ (24 parameters, 4 qubits).

Two strategies were implemented to reduce this error:
\begin{enumerate}
    \item \textbf{Near zero initialization}: Parameters are drawn from $\mathcal{U}(-0.1, 0.1)$ instead of $\mathcal{U}(0, 2\pi)$. At $\theta \approx 0$, approximates the identity by the HWE circuit. The initial state remains close to $|0\ldots0\rangle$ (the Hartree-Fock reference~\parencite{szabo1996modernquantumchemistry} within the Jordan-Wigner mapping~\parencite{peruzzo2014variational}), delaying the particle number violation.

    \item \textbf{Best physical energy tracking}: The optimizer records the lowest energy observed that remains at or above FCI. When the energy crosses below FCI, a single diagnostic message is emitted, and the best physical result is reported instead of the final (unphysical) energy.
\end{enumerate}

It should be stated that this is a previously known limitation of HWE ansatzes, as documented by \textcite{kandala2017hardware} 
in the original HWE VQE demonstration. Due to this, the physically correct alternative, the Unitary Coupled Cluster Singles and Doubles (UCCSD) ansatz~\parencite{romero2018strategiesquantumcomputingmolecular},
was explored during this project but significant integration challenges were encountered with the Qiskit Nature circuit 
library's gate decomposition pipeline. UCCSD still remains architecturally supported within the middleware stack's ansatz tier 
system, kept on as a priority for future updates on the stack.

\subsection{Comparison with Serial Baseline}
Both the distributed stack and its serial baseline use identical, near zero initialization and the same SPSA hyperparameters 
(Table~\ref{table:serial_baseline}), making this an applicable comparison to identify the direct effect of utilizing the MPI distribution. 
Here, the accuracy results are mixed, as the serial baseline achieved a significantly lower error on $LiH$ (2.2 mHa vs. 336.2 mHa) 
and comparable error on $BeH_{2}$ (39.4 mHa for both paths), while the distributed stack performed better on $H_{2}$ (2.4 mHa vs. 11.5 mHa). 
These differences are not caused by an algorithmic advantage of either configuration, due to their identical configurations, but by a stochastic divergence from SPSA 
trajectories due to the use of different execution paths, the distributed \texttt{Allreduce} sum is error 
prone compared to the exact, serial summation, introducing different floating point reduction orders.

The distributed stack's primary contribution is its wall-clock speedup, not its chemical accuracy. Table~\ref{table:speedup} 
shows that with the benchmarked larger molecules that have sufficient Pauli terms to distribute, the stack achieves a $1.11\times$ speedup 
on $BeH_{2}$ and $1.18\times$ speedup on $H_{2}O$, while any smaller molecules will face a net slowdown from the MPI
overhead. This identifies that the parameter initialization strategy has a greater impact on HWE accuracy rather than the execution model.

\subsection{Communication Masking}
The stack architecture supports an asynchronous, non-blocking overlap through the C++ dispatcher's \texttt{std::async} for
QPU submission and the Python path's statevector evaluation during the QPU queue wait times.
The results of the masking metric $M \gg 1$ across all the experiments ran confirms the stack's architectural hypothesis.
Under the assumption that in the simulator path, the classical Pauli term evaluations and statevector construction at a cost of $O(2^n)$ 
will significantly outweigh the cost of MPI communication (two \texttt{MPI\_Allreduce} calls 
per iteration) in the metric $M$. The observed results have implications for the cloud based QPU path, that when the QPU's 
round-trip time ($T_{\text{quant}}$) replaces $T_{\text{comm}}$ as the stack's primary communication bottleneck, the 
classical acceleration workload ($T_{\text{accel}}$) must be substantial enough to mask it. The observed $M$ values 
(ranging between 12--9,000 for the simulator path) confirm that the stack's classical computation is the primary contributor to communication 
overhead by 1--3 orders of magnitude. On the IBM QPU path, $M < 1$ due to the 32--60\,s QPU RTT that vastly exceeds 
the classical workload for a small molecule such as $H_{2}$ (15 Pauli terms). 
For successful QPU masking, the applications of sufficiently complex molecules such as $BeH_2$ or $H_{2}O$ would generate enough classical work to complete during QPU wait times. 

\subsection{IBM QPU: Accuracy Gap and Hardware Noise}
Within its 10 iterations, the IBM QPU achieved $-0.1709\;E_h$ for $H_{2}$, compared to the classical simulator's $-1.1349\;E_h$,
with an error of $0.9663\;E_h$ compared to the correct FCI value. Three main factors explain this error gap. First, 
the QPU run was limited to 10 iterations due to the queue constraints on IBM's open access plan whereas the 
simulator run required 200 iterations for convergence with $H_{2}$ (20x). However, within the 10 iterations there is an identified 
downward energy trend (Table~\ref{table:ibm_iterations}), suggesting that if additional iterations occurred, the quantum hardware would continue to improve the result. 
Secondly, the shot=4096 sampling provides a statistical uncertainty of approximately $1/\sqrt{4096} \approx 0.016\;E_h$ each expectation value estimate, as this produces the unpredictable 
convergence that is shown in Table~\ref{table:ibm_iterations}. Lastly, T-REx of resilience level=1 only handles measurement readout errors without addressing any of the 
gate errors or the now-present issue of decoherence on the physical qubits. To handle these errors, advanced techniques such as Zero Noise Extrapolation (ZNE) 
or Probabilistic Error Cancellation(PEC) would be needed so that hardware noise can be suppressed, as discussed by \textcite{cai2023quantumerrormitigation}.

Despite the QPU's large absolute error, this result is an architectural validation for the end-to-end middleware pipeline when integrated with 
real quantum hardware, including the transpilation pipeline, EstimatorV2 integration, PUB batching, and MPI coordination, which all operate correctly across 10 cloud based QPU iterations without failure.

\subsection{Scaling Limitations and Path to Improvement}
The CPU only scaling results for the molecule $H_{2}O$ show that using MPI distribution on shared-memory hardware will result in a negligible 
benefit, as when compared to MPI $P=1$ it only reached $1.00\times$ at $P=2$ and then regressed to $0.24\times$ at $P=8$. This follows 
Amdahl's Law~\parencite{amdahl1967validity} where the maximum speedup will be bounded by $1/f$ despite of the number of processors utilized. CPU-only MPI
distribution of Pauli terms on the shared memory alone does not produce a meaningful speedup, with the many $O(2^n)$ statevector constructions dominating wall-clock time, whereas the added GPU acceleration does, reducing the per-rank statevector cost for the parallelization to become beneficial.

The GPU-accelerated results (Section~\ref{sec:gpu_results}) confirm this analysis, as with the GPU Aer backend hosted on an 
NVIDIA GTX 1650 Mobile, the larger molecule $H_{2}O$ at $P=8$ achieved a $1.51\times$ speedup over the serial baseline's time and a $1.78\times$ 
speedup over GPU $P=1$. Compare this to the CPU-only path, which regressed down to $0.24\times$ at the same rank. 
The value of GPU path's strong scaling efficiency for $H_{2}O$ at $P=2$ reached $71.1\%$ compared to $50.1\%$ in the CPU only path, 
but most crucially, the GPU path's scaling did not degrade at ranks of $P=4$ or $P=8$, which demonstrates that with GPU acceleration of each rank's 
statevector gate operations reduces the serial fraction sufficiently, allowing for MPI Pauli term distributions to provide continued computational benefits to the stack.

The GPU weak scaling test results further support this finding, with per-iteration time at $P=8$ ranks was $0.239$\,s with GPU vs. the $1.440$\,s CPU only path, providing 
a $6\times$ improvement. This result confirms the statevector construction bottleneck prediction from \textcite{bayraktar2023cuquantum}.

A notable finding is that the cuStateVec backend performed slower than GPU Aer ($242$\,s vs. $161$\,s for $H_{2}O$ at $P=4$), 
resulting due to GTX 1650 Mobile's limited memory bandwidth ($\sim$128 GB/s GDDR6), which is insufficient for the HBM bandwidth that cuStateVec was designed for, as discussed in Section~\ref{sec:methodology}. 

\subsection{Stack Relationship to XACC}
This stack builds on XACC's philosophy~\parencite{mccaskey2018xacc} of handling the QPU as an accelerated coprocessor with a hardware-agnostic and language-agnostic
API, however, it extends their model to address three main limitations in hybrid middleware:
\begin{enumerate}
    \item \textbf{Single process to multi-process}: XACC operates with a single sequential host, but this stack distributes the VQE workload across multiple MPI ranks on one host for a speedup on finding ground state energies of molecules.
    \item \textbf{Synchronous to asynchronous}: XACC's AcceleratorBuffer model is synchronous while this stack implements two complementary, asynchronous paths: the C++ dispatcher uses \texttt{std::async} for non-blocking MPI coordination and local kernel execution, while the Python IBM path uses Qiskit's EstimatorV2 primitive with parallel local statevector evaluation during the cloud QPU wait time.
    \item \textbf{Compiler level to orchestration level}: XACC focuses on circuit compilation and hardware portability but this stack focuses on solving a known hybrid system bottleneck, aiming for accelerated performance of the classical optimization loop surrounding quantum circuit execution.
\end{enumerate}

These two stacks are complementary instead of being competitive, as XACC could easily serve as the circuit compilation 
frontend then to be fed into this project's MPI-distributed optimization backend for a computational speedup with improved chemical accuracy.


\subsubsection{Real World Applications}
The results of this stack have direct implications for real world applications, such as pharmaceutical drug discovery
where molecular ground state energy estimation is one of the main subroutines for predicting the molecular bindings affinities. 
Here, target molecules produce larger Pauli term counts over the 1086 terms from the benchmarked $H_{2}O$, which Section~\ref{sec:results_weak_scaling} 
proved these larger workloads would benefit the most from MPI distribution. Within a material science application, such as a lithium-ion battery
cathode simulation, electronic structure's Hamiltonian calculations map naturally into the stack's Jordan-Wigner pipeline. This stack's validation with 
accessible, lower bound hardware (consumer GPU, IBM free access plan) demonstrates that distributed hybrid VQE is accessible to research groups even without dedicated HPC resources. 